\chapter{Methodology}
\section{The SA-1B Data Engine --- The Flywheel of Generalization}

\subsection{Paradigm Shift: From Manual Annotation to Model-in-the-Loop}
The fundamental challenge in building a foundation model for segmentation is the lack of web-scale segmentation data. Unlike language or image-text pairs abundant on the internet, segmentation masks are not naturally occurring and must be manually created. Traditional approaches would require paying annotators to label billions of masks—economically infeasible and temporally prohibitive.

SAM's solution is elegant but non-obvious: \textbf{leverage the model itself to assist, accelerate, and eventually replace human annotation.} This is the "Data Flywheel" concept—a three-stage iterative process where model improvement unlocks more efficient data collection, which in turn improves the model.

\subsection{Stage 1: Assisted-Manual Annotation (4.3M Masks, 120K Images)}
In the first stage, SAM operates as an \textbf{interactive segmentation assistant}, similar to traditional click-based segmentation tools (e.g., GrabCut, RITM, SimpleClick). The critical insight here is that the model's role is to \textit{accelerate} human labeling, not replace it.

\textbf{The Workflow:}
\begin{enumerate}
    \item Professional annotators view an image in a web-based interface.
    \item They click foreground/background points to indicate an object of interest.
    \item SAM's pre-computed image embedding (from MAE-ViT-H) is queried with these prompts.
    \item The lightweight Prompt Encoder and Mask Decoder run in real-time ($< 50$ms), generating an initial mask prediction.
    \item Annotators refine the mask using pixel-precise brush and eraser tools.
    \item The final mask is saved.
\end{enumerate}

\textbf{Quantitative Dynamics:}
\begin{itemize}
    \item \textbf{Initial annotation time:} 34 seconds per mask (comparable to COCO annotation standards).
    \item \textbf{Model capacity progression:} ViT-B $\rightarrow$ ViT-H during this stage (1 retraining cycle per major improvement).
    \item \textbf{Average masks per image:} $20 \rightarrow 44$ (increased diversity as model improved).
    \item \textbf{Total masks collected:} 4.3M from 120K images.
\end{itemize}

\textbf{Why This Matters:}
Unlike fully automatic approaches that can only label what the model "already knows," human annotators introduce \textit{novel object categories, rare poses, and edge cases} that force the model to generalize. Critically, annotators were \textbf{not constrained to semantic categories}—they could label "stuff" (amorphous regions like sky, grass) and "things" (discrete objects like people, cars) uniformly, maximizing the diversity of training signals.

\textbf{Efficiency Gain Analysis:}
The $34 \rightarrow 14$ second progression represents a \textbf{6.5$\times$ speedup over COCO standards} (34 seconds here vs. 14 seconds for bounding boxes). This exponential improvement in annotation velocity is crucial for economic feasibility at billion-scale.

\subsection{Stage 2: Semi-Automatic Annotation (5.9M Additional Masks, 180K Images)}
In Stage 2, the goal shifts explicitly: \textbf{increase mask diversity}. The intuition is that Stage 1 naturally biases toward prominent, easy-to-label objects (due to annotators' tendency to prioritize high-salience items). Stage 2 forces the model to learn harder, less obvious objects.

\textbf{The Semi-Automatic Protocol:}
\begin{enumerate}
    \item Train a generic object detector on all masks from Stage 1. This detector outputs bounding boxes regardless of semantic class (object-agnostic proposal generation).
    \item SAM automatically segments high-confidence regions identified by this detector, generating candidate masks.
    \item Annotators are presented with images pre-filled with these automatic masks.
    \item Annotators' task: Label only the \textit{unannotated} objects (those not covered by the automatic masks).
    \item This forces annotators to focus on smaller, more occluded, or semantically ambiguous objects.
\end{enumerate}

\textbf{Quantitative Dynamics:}
\begin{itemize}
    \item \textbf{Annotation time (excluding automatic masks):} $\sim$34 seconds per mask (unchanged; challenging objects require the same effort).
    \item \textbf{Average masks per image:} $44 \rightarrow 72$ (including automatic masks; diversity increased substantially).
    \item \textbf{New masks collected:} 5.9M, raising the cumulative total to 10.2M masks.
    \item \textbf{Retraining cycles:} 5 full model retraining iterations.
\end{itemize}

\textbf{Why This Engineering is Critical:}
The semi-automatic stage is a \textbf{curriculum learning} strategy. By automatically handling easy cases and forcing annotators to label hard cases, the model learns a more balanced representation. This is empirically validated: the model's ability to handle edge cases improved significantly from Stage 1 to Stage 2, as evidenced by the diversity metrics.

\textbf{Annotation Bias Mitigation:}
Without Stage 2, the model would develop a strong bias toward dominant objects (people, vehicles, large structures). The semi-automatic strategy explicitly breaks this bias by requiring exhaustive annotation of all objects, regardless of prominence.

\subsection{Stage 3: Fully Automatic Annotation (1.1B Masks, 11M Images)}
This is the stage where the data flywheel reaches critical velocity. With sufficient diverse masks from Stages 1 and 2, the model is now capable enough to generate masks \textbf{without any human intervention}.

\subsubsection{The $32 \times 32$ Grid Prompting Strategy}
The core insight is deceptively simple: \textbf{systematically prompt the model with a regular spatial grid of points, covering the entire image.}

\textbf{Mathematical Formulation:}
\begin{itemize}
    \item \textbf{Grid resolution:} $32 \times 32 = 1024$ grid points per image.
    \item Each point is passed as a prompt to the Prompt Encoder.
    \item For each point, SAM predicts \textbf{three masks simultaneously} (owing to its ambiguity-aware design).
    \item \textbf{Theoretical maximum:} $1024 \text{ points} \times 3 \text{ masks per point} = 3072 \text{ masks per image}$.
\end{itemize}

\textbf{Why This Grounding Works:}
A regular grid ensures comprehensive spatial coverage. If an object's boundary passes through or near a grid point, the model will capture it from that point. Objects of varying sizes are covered because the grid is at a fixed, moderate resolution. This is fundamentally different from random point sampling, which can miss objects entirely.

\subsubsection{Multi-Scale Masking and Subpart Prediction}
A critical architectural feature of SAM (and one of the reasons Stage 3 is feasible) is its \textbf{ambiguity-aware decoder} that outputs multiple masks per prompt.

\textbf{The Output Token Design:}
\begin{itemize}
    \item The Mask Decoder predicts three masks from a single prompt, ranked by IoU score.
    \item \textbf{Whole object:} The complete, semantically coherent object.
    \item \textbf{Part:} A prominent sub-component (e.g., a wheel on a car, a face on a person).
    \item \textbf{Sub-part:} Fine-grained components (e.g., an eye within a face, a spoke within a wheel).
\end{itemize}

\textbf{Why Three Outputs?}
Empirical analysis during SAM's development showed that image regions typically have at most three levels of semantic hierarchy (whole $>$ part $>$ sub-part). Predicting three masks covers all common ambiguity scenarios without excessive computation.

\textbf{Mathematical Formulation:}
During training, SAM employs \textbf{minimum loss} over the three mask outputs:
\[
\mathcal{L}_{\text{total}} = \min(\mathcal{L}_{i=1}^3) \quad \text{for masks } m_1, m_2, m_3
\]
This ensures that at least one of the three masks aligns with ground truth, forcing the model to learn diverse interpretations of ambiguous prompts.

\subsubsection{Filtering, Deduplication, and Non-Maximum Suppression}
\textbf{The Critical Question:} If a $32 \times 32$ grid yields up to 3072 theoretical masks, why do we only get $\sim$100 high-quality masks per image?

\textbf{Answer:} Aggressive filtering at multiple stages.

\textbf{Stage 3.A: IoU Scoring and Confidence Thresholding}
\begin{itemize}
    \item For each of the three masks predicted from a grid point, SAM outputs a \textbf{predicted IoU score} (estimated mask quality).
    \item IoU prediction is a separate output head in the Mask Decoder: a learned estimate of the mask's alignment with the true object boundary.
    \item \textbf{Threshold:} Only masks with \textbf{predicted IoU $>$ 0.88} are retained.
    \item \textbf{Rationale:} High IoU threshold ensures only high-quality masks enter the deduplication pipeline.
    \item \textbf{Expected reduction:} From 3072 theoretical outputs, this threshold alone eliminates $\sim$80\% of predictions (keeping only predictions the model is highly confident about).
\end{itemize}

\textbf{Stage 3.B: Stability Filtering}
\begin{itemize}
    \item A mask is \textbf{stable} if small perturbations in the model's confidence threshold (e.g., changing the foreground probability threshold from 0.5 to $0.5 \pm \epsilon$) produce similar masks.
    \item \textbf{Operationally:} Threshold the mask's probability map at both 0.5 and 0.5, and compute the IoU between the two resulting binary masks.
    \item If this "stability IoU" is high ($> 0.95$), the mask is kept; otherwise, it is discarded.
    \item \textbf{Why This Matters:} Masks that are brittle (sensitive to small threshold changes) are often segmentation errors or include spurious boundaries. Stability filtering removes low-confidence predictions and keeps robust ones.
    \item \textbf{Further reduction:} $\sim$50\% of remaining masks fail stability checks.
\end{itemize}

\textbf{Stage 3.C: Non-Maximum Suppression (NMS) for Mask Deduplication}
After IoU and stability filtering, many remaining masks are \textbf{duplicates or near-duplicates} (same object segmented from slightly different grid points, or different granularities of the same object).

\textbf{Traditional NMS for Bounding Boxes:}
Classical NMS (used in object detection) operates on bounding boxes and IoU:
1. Sort boxes by confidence score.
2. Select the highest-confidence box.
3. Remove all boxes with IoU $>$ threshold to this box.
4. Repeat on remaining boxes.

\textbf{Mask-Based NMS in SAM:}
The SAM paper adapts NMS for mask deduplication:
1. Sort masks by predicted IoU score.
2. Select the highest-scoring mask $m_1$.
3. For all remaining masks $m_i$:
   \begin{itemize}
       \item Compute $\text{IoU}(m_1, m_i) = \frac{|m_1 \cap m_i|}{|m_1 \cup m_i|}$ (Jaccard index).
       \item If $\text{IoU}(m_1, m_i) > \text{thresh}_{\text{IoU}}$ (typically 0.7), discard $m_i$ (it is a near-duplicate).
       \item If $\text{IoU}(m_1, m_i) \leq \text{thresh}_{\text{IoU}}$, retain $m_i$ (it is a distinct object).
   \end{itemize}
4. Repeat on remaining unprocessed masks.

\textbf{Result:} NMS eliminates the redundant foreground predictions while preserving distinct objects, even when they overlap partially.

\textbf{Scale Reduction Summary:}
\begin{itemize}
    \item \textbf{Start:} 3072 theoretical masks ($32 \times 32$ grid $\times$ 3 outputs per grid cell).
    \item \textbf{After IoU filtering ($>0.88$):} $\sim$600 masks (20\% remain).
    \item \textbf{After stability filtering:} $\sim$300 masks (50\% of IoU-filtered).
    \item \textbf{After NMS ($\text{thresh\_IoU} = 0.7$):} $\mathbf{\sim100 \text{ masks}}$ (final).
\end{itemize}

\textbf{Final Calculation:}
\[
\text{Total SA-1B Masks} = 11\text{M images} \times 100 \text{ masks/image} = 1.1 \text{ B masks}
\]

\subsubsection{Multi-Scale Processing: Zoomed-in Image Crops}
To further improve small object segmentation (which is notoriously challenging), SAM processes \textbf{multiple overlapping zoomed-in crops} of each image.

\textbf{The Zooming Strategy:}
\begin{itemize}
    \item Original image is processed at full resolution (1500 pixels shortest side after downsampling).
    \item Additionally, crop and zoom into regions systematically (e.g., $2\times, 4\times$ magnification).
    \item The $32 \times 32$ grid prompting is re-applied to these zoomed crops.
    \item Additional fine-grained masks are generated for small objects that might be missed in the full-resolution pass.
    \item These zoomed-crop masks are appended to the final set (before final NMS).
\end{itemize}

\textbf{Why This Is Crucial:}
Small objects (relative size $< 1\%$ of image) have very few pixels. A standard NMS threshold might inadvertently merge them with larger neighbors. By processing zoomed crops, SAM explicitly allocates model capacity to these hard-to-segment regions.